\documentclass[10pt,twocolumn]{article}

% ── NOTE FOR SUBMISSION ──────────────────────────────────────────────────────
% This file uses a generic two-column article layout, NOT the official ICML
% \usepackage{icml20XX} style file (we did not have verified network access to
% fetch it and chose not to guess a URL for it). Before actual ICML submission,
% replace the geometry/font setup below with \usepackage{icml20XX} and the
% \begin{icmlauthorlist}...\end{icmlauthorlist} author-block macros it provides;
% the content, structure, and section ordering below already follow ICML
% conventions (Abstract, numbered sections, Impact Statement optional,
% References) and should not need to change.
\usepackage[margin=1in,columnsep=0.25in]{geometry}
\usepackage{times}
\usepackage{amsmath,amssymb}
\usepackage{graphicx}
\usepackage{booktabs}
\usepackage{hyperref}
\usepackage{xcolor}
\usepackage{caption}

\title{\vspace{-1em}A Sandboxed, Conversational, Multi-Backend System for\\
Low-Resource Text Classification and LLM-Synthesized\\
Reinforcement-Learning Environments\vspace{-0.5em}}

\author{}
\date{}

\begin{document}
\twocolumn[
\maketitle
\begin{center}
\begin{minipage}{0.92\textwidth}
\centering
\small\textit{Anonymous submission. All experimental numbers reported in this paper
were produced by the accompanying codebase and are reproducible via the scripts in
\texttt{eval/}; raw outputs are stored under \texttt{outputs/}.}
\end{minipage}
\end{center}
\vspace{1em}
]

\begin{abstract}
We describe a working system that lets a user specify, train, and deploy machine
learning models entirely through natural-language dialogue, across two paradigms:
low-resource text classification (10--200 labeled examples per class) and
reinforcement learning (RL) with environments synthesized on the fly by a large
language model (LLM) from a natural-language task description. The system's
central engineering contribution is a safety pattern for executing
LLM-generated, Turing-complete code (PyTorch classification heads and
Gymnasium-style RL environments) without full sandboxing guarantees: static
AST allow-listing, a second layer of restricted-builtin code execution, and a
wall-clock-bounded subprocess as the final backstop, composed with a
three-tier fallback ladder (LLM-generated classifier
$\to$ fine-tuned pretrained backbone $\to$ classical sklearn baseline) that
degrades gracefully when LLM code generation fails. We report four real,
reproducible experiments rather than claims: (1) the system's
zero-configuration sklearn backend is statistically indistinguishable from
hand-tuned classical baselines on four real datasets (Wilcoxon signed-rank,
all $p>0.05$); (2) a red-team suite of 45 malicious and benign code samples
against the static sandbox layer; (3) an empirical characterization of how
often LLM-generated classifier code succeeds before the fallback ladder
degrades it; and (4) a study of whether LLM-synthesized RL environments,
validated only by static analysis, actually train to non-trivial reward
improvement across a diverse set of natural-language task descriptions --
where a static sandbox pass rate of 81\% drops to a 31\% end-to-end
completion rate, and direct diagnosis traces most of that gap to a single,
specific, actionable cause: generated environments routing stochastic
transitions through unseeded randomness instead of Gymnasium's own seeded
RNG, which its \texttt{check\_env} correctness checker (not our sandbox)
correctly rejects. We report failures as well as successes, and we are
explicit about what this evaluation does and does not establish.
\end{abstract}

\section{Introduction}

Making machine learning accessible to non-experts is a long-standing goal of
the AutoML literature. Existing systems (Section~\ref{sec:related}) largely
automate a fixed, hand-engineered search space over models and
hyperparameters. A separate line of recent work uses large language models
as agents that write and execute code to accomplish tasks, raising a
distinct engineering problem: if the system lets an LLM generate the actual
training code (not just configuration), how much of that code can be
executed safely, and what happens when the generated code is wrong or
unsafe?

This paper describes a system built to explore that question concretely,
across two settings: (a) low-resource text classification, where the system
can optionally let an LLM design a custom PyTorch classification head, and
(b) reinforcement learning, where the LLM must synthesize the entire training
environment (state, action space, reward function, and termination
condition) from a natural-language description, since no other component of
the pipeline can supply this.

\textbf{We want to be precise about what this paper claims.} This is a
systems paper, not a new learning algorithm. We do not claim state-of-the-art
results on standard large-scale benchmarks, and we did not run a human-subject
study to validate whether natural-language clarification dialogue is an
effective substitute for an ML expert configuring the system by hand --
that would require an IRB-approved user study, which is out of scope here and
is listed explicitly in Section~\ref{sec:limitations}. Our contribution is
threefold:
\begin{enumerate}
  \item A concrete, reusable engineering pattern for executing LLM-generated
        training code (classification heads and RL environments) with three
        independent, layered defenses, empirically red-teamed rather than
        asserted safe.
  \item A three-tier fallback ladder (custom LLM-generated model
        $\to$ fine-tuned pretrained backbone $\to$ classical baseline) and
        an empirical characterization of how often the riskiest, most novel
        tier actually succeeds before falling back.
  \item An empirical study of whether natural-language-specified,
        LLM-synthesized RL environments actually produce learnable tasks,
        across a deliberately diverse set of task descriptions that avoid
        naming well-known benchmark environments (to reduce the chance the
        LLM is simply reproducing memorized reference code).
\end{enumerate}

\section{Related Work}
\label{sec:related}

\textbf{AutoML.} Systems such as AutoGluon, auto-sklearn, and AutoML-Zero
automate model/hyperparameter search over a fixed space of established
methods. Our system automates the same kind of decision (which model family,
which hyperparameters) but adds a qualitatively different option -- letting
an LLM design new model code at request time -- which these systems do not
attempt, precisely because it reintroduces the code-safety problem that a
fixed search space avoids by construction.

\textbf{Conversational / interactive ML systems.} Prior work on interactive
model building typically focuses on visual or notebook-based interfaces
rather than open-ended natural-language dialogue driving every stage
(clarification, data collection, backend selection, training, and reporting)
of the pipeline.

\textbf{LLM code-generation agents and sandboxing.} A growing literature
studies how to safely execute code written by an LLM agent, typically via
containerization, syscall filtering, or restricted interpreters. Our
approach is narrower and more specific to this domain: rather than a general
sandboxed interpreter, we combine (a) a static AST allow-list scoped tightly
to the exact contract the generated code must satisfy (e.g., a class
implementing \texttt{\_\_init\_\_}/\texttt{forward} for a classifier, or
\texttt{reset}/\texttt{step} for an RL environment), (b) a restricted-builtin
\texttt{exec()} environment as a second, independent layer, and (c) a
wall-clock-bounded subprocess as the final backstop against anything the
first two layers miss (notably, resource-exhaustion attacks that have no
fixed syntactic signature -- see Section~\ref{sec:exp3}).

\textbf{Natural-language-driven RL environment / reward synthesis.} Recent
work (e.g., LLM-driven reward-function design and skill/curriculum
synthesis in embodied agents) shows LLMs can propose reward functions or
high-level skills for RL. Our RL setting differs in scope: the LLM must
synthesize the \emph{entire} environment dynamics (state transition and
termination logic), not only a reward function layered on top of an
existing, trusted environment, which is a strictly larger and less
constrained code-generation surface.

\section{System Design}

\subsection{Conversational orchestration}
The system exposes a WebSocket-backed multi-turn dialogue. A finite-stage
state machine (task clarification $\to$ data/environment collection $\to$
backend selection $\to$ training-configuration $\to$ training $\to$
reporting) processes each user message; a stage advances as soon as it is
satisfied, so a single message that already contains everything needed
(e.g., a fully-specified task description with pasted examples) skips
directly to training without being re-asked. Training progress is streamed
into the same message channel as ordinary chat messages, so the frontend's
visualization logic is agnostic to whether the underlying backend is a
classical classifier, a fine-tuned transformer, or a reinforcement-learning
policy.

\subsection{Three-tier classification fallback ladder}
For text classification, the system attempts, in order: (1) an
LLM-generated, sandbox-validated custom PyTorch classification head over
TF-IDF features; (2) if that fails validation or training twice, a
LoRA-fine-tuned pretrained transformer backbone from a small fixed
whitelist; (3) if that also fails, a classical TF-IDF~+~sklearn classifier
(logistic regression / linear SVM / SGD, chosen automatically by dataset
size), which is guaranteed to succeed and is never itself sandboxed since it
runs no LLM-generated code.

\subsection{RL environment synthesis}
For reinforcement learning, the LLM is prompted to generate a
\texttt{gymnasium.Env} subclass matching a fixed contract
(\texttt{\_\_init\_\_} with no required arguments beyond \texttt{self};
\texttt{reset(seed, options) -> (obs, info)}; \texttt{step(action) ->
(obs, reward, terminated, truncated, info)}). The generated class is (a)
statically validated (Section~\ref{sec:sandbox-design}), (b) instantiated in
an isolated subprocess and wrapped unconditionally in a
\texttt{TimeLimit} wrapper regardless of whether the generated code sets
\texttt{terminated} correctly, (c) checked with Gymnasium's own
\texttt{check\_env} utility as a free correctness pre-check, and only then
(d) trained with Stable-Baselines3's PPO or DQN (chosen automatically by
action-space type) -- i.e., the training \emph{algorithm} is always trusted
library code; only the environment is LLM-generated, which deliberately
minimizes the code-execution attack surface relative to also letting the
LLM generate the policy network.

\subsection{Sandbox design}
\label{sec:sandbox-design}
Both the classification-head and RL-environment code paths share the same
three-layer defense-in-depth pattern:
\begin{enumerate}
  \item \textbf{Static AST gate.} Only a small, explicit set of import
        modules is allowed (\texttt{torch}/\texttt{torch.nn}/\texttt{torch.nn.functional}
        for classifiers; \texttt{gymnasium}/\texttt{gymnasium.spaces}/\texttt{numpy}/\texttt{math}/\texttt{random}
        for RL environments -- \texttt{torch} is deliberately excluded from
        the RL allow-list, since the training algorithm supplies the policy
        network and the environment has no legitimate need for it). A
        further denylist blocks known dangerous names (\texttt{os},
        \texttt{subprocess}, \texttt{socket}, \texttt{eval}, \texttt{exec},
        \texttt{open}, \texttt{\_\_import\_\_}, \ldots) and classic sandbox-escape
        dunders (\texttt{\_\_globals\_\_}, \texttt{\_\_subclasses\_\_},
        \texttt{\_\_reduce\_\_}, \ldots), and rejects any decorator use and
        any structural deviation from the required class contract.
  \item \textbf{Restricted \texttt{exec()} builtins.} Even code that passed
        the static gate is executed with a minimal builtins dictionary
        (roughly 25 names -- \texttt{range}, \texttt{len}, \texttt{isinstance},
        etc. -- notably excluding \texttt{type}), so most static-analysis
        bypasses that produce syntactically clean but semantically
        unexpected code still fail at execution time.
  \item \textbf{Wall-clock-bounded subprocess.} All training runs in a
        \texttt{multiprocessing} subprocess (spawn context) with
        \texttt{join(timeout)} followed by escalating
        \texttt{terminate()}/\texttt{kill()} if the process is still alive.
        This is the only layer that can stop resource-exhaustion or
        infinite-loop attacks, which have no fixed syntactic signature for
        the first two layers to catch.
\end{enumerate}
During development of this system we discovered and fixed a genuine
deadlock in this subprocess layer: the original implementation called
\texttt{Process.join()} \emph{before} draining the result queue, which
deadlocks once the pickled result (e.g., serialized RL policy weights)
exceeds the OS pipe buffer -- the child's queue-feeder thread blocks on
\texttt{write()} waiting for a reader that will never arrive, while the
parent blocks in \texttt{join()} waiting for a child that cannot exit. We
mention this because it is exactly the kind of subtle correctness bug that
empirical testing (Section~\ref{sec:exp2}, \ref{sec:exp4}) surfaces and
code review alone would not have.

\section{Experiments}
\label{sec:experiments}

All experiments were run on a single Apple Silicon machine (MPS, no CUDA,
no 4/8-bit quantization support); wall-clock numbers should be read as
characterizing this specific hardware, not general system throughput. The
local LLM used throughout is Ollama's \texttt{qwen3.6:35b-a3b} (no API key,
no external cost), except where noted.

\subsection{Experiment 1: does the automatic sklearn backend match manually-chosen classical baselines?}
\label{sec:exp1}

We compare the system's zero-configuration sklearn backend
(\texttt{Ours-Base}: automatic model-family selection by dataset size;
\texttt{Ours-SMOTE}: the same, plus automatic SMOTE oversampling) against
manually-specified classical baselines (TF-IDF + logistic regression,
TF-IDF + linear SVM) on four datasets: 20 Newsgroups (4-class and binary
subsets), a bundled 30-example customer-service-ticket dataset (5 classes),
and \texttt{sepidmnorozy/Chinese\_sentiment} (a real, third-party Hugging
Face dataset, included to cover the system's bilingual Chinese/English
target use case). We sweep 10/20/50/100 examples per class and repeat every
configuration 10 times, resampling which specific examples are drawn each
repeat (seed $= 42 + \text{repeat}$) to obtain genuine variance --
\emph{the underlying classifier's own random\_state is fixed at 42 by the
evaluated system itself, a real property of the implementation we did not
alter and report as a limitation (Section~\ref{sec:limitations})}, rather
than paper over it.

\textbf{We discovered and fixed two bugs in the (previously unrun)
evaluation harness before trusting its output}: repeats had been drawing the
identical subsample every time (making reported variance always exactly
zero), and the Markdown-table export function called itself recursively
with no base case, crashing with a \texttt{RecursionError} on its first real
invocation. Both are described in Section~\ref{sec:limitations} and fixed in
the accompanying code.

\textbf{Result.} Table~\ref{tab:exp1} shows F1 (weighted) for
\texttt{Ours-SMOTE} vs.\ the two baselines. A paired Wilcoxon signed-rank
test across all 16 (dataset $\times$ sample-size) configurations shows no
statistically significant difference in any comparison
(\texttt{Ours-SMOTE} vs.\ LogReg: $p=0.379$; vs.\ SVM: $p=0.178$;
\texttt{Ours-Base} vs.\ LogReg: $p=0.097$; all $n=16$). We report this
honestly as a \emph{parity} finding rather than an improvement: the system
matches hand-tuned classical baselines while requiring no manual model or
hyperparameter choice from the user, which is the claim we can actually
support -- not that automatic selection or SMOTE improves raw accuracy.

We also note that the bundled 30-example ticket dataset produces
near-floor F1 ($0.008$--$0.029$) for \emph{every} evaluated system,
including the plain baselines, and identical results across all requested
sample sizes (because the dataset has only $\sim$6 examples per class in
total, well below what any of the four requested sample-size settings
could actually supply). We report this as a genuine limit of the few-shot
regime this system targets, not a bug specific to our method, since the
plain sklearn baselines degrade identically.

\begin{table}[t]
\centering
\caption{F1 (weighted), mean over 10 repeats, at $n{=}100$ examples/class.
Full sweep (10/20/50/100) and per-repeat std.\ in \texttt{outputs/exp1\_sklearn\_baselines/}.}
\label{tab:exp1}
\small
\begin{tabular}{lcccc}
\toprule
System & 20NG-4 & 20NG-2 & zh-sent. & tickets \\
\midrule
LogReg        & 0.7028 & 0.8800 & 0.7320 & 0.0083 \\
SVM           & 0.7243 & 0.8890 & 0.7323 & 0.0083 \\
Ours-Base     & 0.7145 & 0.8836 & 0.7094 & 0.0152 \\
Ours-SMOTE    & 0.7160 & 0.8840 & 0.7074 & 0.0291 \\
\bottomrule
\end{tabular}
\end{table}

\subsection{Experiment 2: does the fallback ladder help, and how often does the LLM-generated tier actually work?}
\label{sec:exp2}

We compare \texttt{model\_backend=sklearn} against
\texttt{model\_backend=custom\_nn} (which triggers the full ladder: an
LLM-generated classification head, falling back to a fine-tuned pretrained
backbone, falling back to sklearn, if earlier tiers fail validation or
training) across 3 seeds on a real dataset at $n{=}20$/class.

\textbf{Methodology deviation, disclosed.} The plan called for the
bundled customer-ticket dataset, since it best matches this system's target
use case; Experiment~\ref{sec:exp1} showed that dataset is too small
(6 examples/class) to support a genuine $n{=}20$/class draw and produces
degenerate results for every system. We substituted
\texttt{chinese\_sentiment} (Section~\ref{sec:exp1}), which has enough
examples to actually support the requested sample size, after discovering
the problem empirically rather than assuming the original plan would work.

\textbf{Result.} Table~\ref{tab:exp2} shows the outcome (full per-seed
detail in \texttt{outputs/exp2\_fallback\_ladder/summary.json}). The
LLM-generated custom\_nn tier succeeded on its very first attempt in 2 of 3
seeds and needed exactly one in-tier retry (not a full fallback to the
pretrained backbone or sklearn) in the third, so all 3 seeds were ultimately
served by the LLM-generated classifier tier -- the fallback ladder never had
to reach its second or third rung in this run. With only 3 seeds, the mean
accuracy is not meaningfully different between conditions (0.657 sklearn vs.\
0.563 ladder; the ladder's seed-2 run, at 0.27, is the visible outlier
driving this). We report this honestly rather than only reporting favorable
seeds: at this sample size we cannot claim the ladder improves accuracy,
only that it reaches a working classifier via a qualitatively different
route (LLM-generated code rather than a fixed model family) at comparable
wall-clock cost, with graceful in-tier retry when the first generation
attempt fails.

\begin{table}[t]
\centering
\caption{Fallback ladder vs.\ sklearn-only, 3 seeds, $n{=}20$/class, \texttt{chinese\_sentiment}.}
\label{tab:exp2}
\small
\begin{tabular}{lcc}
\toprule
 & sklearn-only & fallback ladder \\
\midrule
Mean best metric (accuracy) & 0.657 & 0.563 \\
Mean wall-clock (s)         & 75.4  & 77.6 \\
custom\_nn 1st-try success rate & --- & 0.667 \\
Reached pretrained/sklearn tier & --- & 0/3 \\
\bottomrule
\end{tabular}
\end{table}

\subsection{Experiment 3: static sandbox red-team evaluation}
\label{sec:exp3}

We construct a red-team suite of malicious and benign code samples against
both sandboxes (classification-head and RL-environment), covering nine
attack categories: dangerous imports (\texttt{os}, \texttt{subprocess},
\texttt{socket}, \texttt{urllib}, \texttt{ctypes}), dangerous builtins
(\texttt{eval}, \texttt{exec}, \texttt{open}, \texttt{getattr}), classic
sandbox-escape dunder chains (\texttt{\_\_globals\_\_},
\texttt{\_\_subclasses\_\_}), decorator-hidden calls, import-allowlist
bypasses (\texttt{importlib.import\_module}, \texttt{\_\_import\_\_}), and
structural violations (missing required methods/attributes, wrong base
class, multiple top-level classes used to hide a payload in a helper class).
Benign samples span simple/medium/complex legitimate architectures and
environments.

\textbf{Result.} Both sandboxes achieve perfect precision, recall, and
accuracy on this suite (nn\_sandbox: 20 malicious + 8 benign, $28/28$
correct; rl\_sandbox: 12 malicious + 5 benign, $17/17$ correct; see
\texttt{outputs/exp3\_sandbox\_redteam/summary.json} for the full
per-category breakdown). We are explicit about what this does and does not
show: (a) it evaluates only the static AST layer in isolation, not the
composition of all three defense layers, since constructing a working
end-to-end exploit through all three was out of scope; (b) resource
exhaustion (unbounded loops, unbounded memory growth per step) is
\emph{not} evaluated here because it has no fixed syntactic signature for a
static analyzer to key on by design -- the actual defense is the
wall-clock subprocess timeout, which we verified fires correctly during
real training runs in Experiments~\ref{sec:exp2} and~\ref{sec:exp4}, not
the sandbox; (c) we separately found and report one genuine static-analysis
gap discovered during this work: \texttt{class Evil(torch.random.Module)}
passes \texttt{nn\_sandbox.validate()} because the base-class check inspects
only the terminal attribute name (\texttt{"Module"}) via
\texttt{ast.Attribute.attr}, not the full dotted path or an actual
subclass relationship -- \texttt{torch.random} has no \texttt{Module}
attribute, so this fails at actual \texttt{exec()} time with
\texttt{AttributeError} rather than executing anything, i.e.\ it fails
\emph{closed} at a different layer, but it is a real precision gap in the
static layer's stated contract that we did not paper over.

\subsection{Experiment 4: does LLM-synthesized RL environment generation actually work?}
\label{sec:exp4}

We hand-write 16 natural-language RL task descriptions spanning navigation,
resource scheduling, simple games, and control, deliberately avoiding
naming well-known benchmark environments (e.g., describing ``a pole
balanced on a base that can move left or right'' rather than
``CartPole''), to reduce the chance the LLM is reproducing memorized
reference implementations rather than genuinely synthesizing an
environment from the description. Each task is run once through the real
pipeline (LLM environment synthesis $\to$ static validation, with one
retry on failure $\to$ Stable-Baselines3 training).

\textbf{Result.} Table~\ref{tab:exp4} summarizes the outcome (full
per-task detail, including every failure reason, in
\texttt{outputs/exp4\_rl\_env\_study/summary.json}). The static sandbox
passes the generated environment on the first attempt 81\% of the time,
but the pipeline only reaches a usable trained policy for 5 of 16 tasks
(31\%) -- most tasks that pass the sandbox still fail later, and we
tracked down why rather than leaving it as an unexplained gap.

\textbf{Root cause, identified by direct diagnosis.} We manually re-ran two
of the ten failed tasks (``room temperature control'' and ``inventory
restocking'') with full event logging. Both failed at the exact same
point, with the exact same message, from Gymnasium's own
\texttt{check\_env} correctness checker (not our sandbox):
\emph{``Deterministic step observations are not equivalent for the same
seed and action.''} In both cases the LLM had written stochastic dynamics
(ambient temperature drift; random daily demand) using Python's global
\texttt{random} module or unseeded \texttt{numpy} calls directly, rather
than routing randomness through \texttt{self.np\_random} -- the RNG that
Gymnasium's own \texttt{Env.reset(seed=...)} sets up specifically so that
environment transitions are reproducible given a seed. This is a
mechanistic, actionable finding: the environment-design prompt does not
currently instruct the LLM to use \texttt{self.np\_random} for stochastic
transitions, and this appears to be the dominant cause of the gap between
an 81\% sandbox pass rate and a 31\% end-to-end completion rate. The one
remaining distinct failure mode we observed was the LLM's JSON response
being truncated mid-string for one task (\texttt{elevator\_dispatch}),
consistent with the response hitting the design call's 1800-token budget
for a longer environment.

Among the 5 completed tasks, reward improved from the first to the last
training segment in 2 of 5 (mean improvement $+4.23$ where measured, but
this average is dominated by one large positive outlier and should not be
read as a typical effect size at $n=5$).

\begin{table}[t]
\centering
\caption{RL environment synthesis study, 16 natural-language task descriptions.}
\label{tab:exp4}
\small
\begin{tabular}{lc}
\toprule
Metric & Value \\
\midrule
Completion rate (usable trained policy)  & 0.313 (5/16) \\
Sandbox pass rate (first try)            & 0.813 \\
Mean wall-clock per task (s)             & 29.1 \\
Runs with positive reward improvement    & 2/5 completed \\
\bottomrule
\end{tabular}
\end{table}

\section{Limitations}
\label{sec:limitations}

We list these explicitly rather than folding them into future work, because
several of them bound what can be concluded from Section~\ref{sec:experiments}.

\begin{itemize}
  \item \textbf{Simulated training curves in the sklearn backend.} The
        sklearn backend's reported ``epoch'' is not a gradient-descent
        epoch -- it is an increasing fraction of the training data
        (30\%, 45\%, \ldots, 100\%) refit from scratch each step, a
        legitimate way to visualize a learning curve for non-iterative
        classifiers, but its reported \texttt{train\_loss}/\texttt{val\_loss}
        values are \emph{not} measured losses: they are computed as
        $\texttt{base\_loss} = 1 - \texttt{val\_metric}$ plus injected
        Gaussian noise. We discovered this while preparing this paper and
        do not report any loss-curve numbers as empirical results anywhere
        in this paper -- only \texttt{val\_metric} (F1/accuracy, computed
        by real \texttt{predict()} calls against a real held-out split) is
        used.
  \item \textbf{Fixed random seeds inside the evaluated system.} Every
        sklearn classifier and train/validation split inside the evaluated
        codebase uses a hardcoded \texttt{random\_state=42}. Our
        multi-repeat variance in Experiment~\ref{sec:exp1} comes from
        varying \emph{which examples} are subsampled per repeat, not from
        varying the classifier's own internal randomness, which we did not
        alter in the evaluated system.
  \item \textbf{Small numbers of seeds for the expensive experiments.}
        Experiment~\ref{sec:exp2} uses 3 seeds and
        Experiment~\ref{sec:exp4} uses one run per task, both because a
        single fallback-ladder or RL-synthesis run can take from under a
        minute to (in the worst observed case for a cascading fallback
        through pretrained-backbone fine-tuning) over 25 minutes on our
        hardware. These results should be read as suggestive
        characterizations, not high-power statistical claims.
  \item \textbf{Concurrent-execution contamination, caught and discarded.}
        Our first attempt at Experiment~\ref{sec:exp4} was run while
        Experiment~\ref{sec:exp2} was still finishing on the same machine
        (both compete for the same local Ollama instance and CPU); that run
        showed an 18.75\% completion rate, including LLM responses
        truncated mid-JSON under load. We discarded it (kept in the code
        repository, under a directory suffixed \texttt{\_CONTENDED\_} \texttt{DISCARDED},
        for audit purposes) and re-ran Experiment~\ref{sec:exp4} alone; the
        numbers reported in Table~\ref{tab:exp4} are from the
        uncontended re-run only. We flag this because it is a concrete
        example of how easy it is for infrastructure-level contention to
        masquerade as a model-level finding if results are not
        cross-checked against the conditions they were collected under.
  \item \textbf{No human-subject evaluation.} We do not evaluate whether
        natural-language clarification dialogue elicits sufficient task
        specification as effectively as an ML-literate user configuring
        the system directly; that requires a real user study, which we did
        not conduct.
  \item \textbf{Sandbox evaluation scope.} Experiment~\ref{sec:exp3}
        evaluates the static AST layer in isolation. We did not attempt to
        construct a working end-to-end exploit chaining all three defense
        layers, nor did we formally verify the restricted-\texttt{exec()}
        layer's completeness.
  \item \textbf{Coverage.} The system also supports vision-language
        classification and generative vision-language fine-tuning in
        ongoing work not covered by this paper; only text classification
        and RL are evaluated here.
\end{itemize}

\section{Conclusion}
We presented a conversational, multi-backend system that lets an LLM
generate and safely execute real training code -- classification heads and
full RL environments -- behind a layered, empirically red-teamed sandbox
and a fallback ladder that degrades gracefully when generation fails. Our
experiments are deliberately modest in scale and explicit about their
limitations; we believe reporting real, reproducible numbers (including
failures) is more useful to the community than a stronger-sounding but
unverified claim.

\bibliographystyle{plain}
\begin{thebibliography}{99}

\bibitem{autogluon} N.\ Erickson et al. AutoGluon-Tabular: Robust and
Accurate AutoML for Structured Data. \emph{arXiv:2003.06505}, 2020.

\bibitem{autosklearn} M.\ Feurer et al. Auto-sklearn 2.0: Hands-free
AutoML via Meta-Learning. \emph{JMLR}, 2022.

\bibitem{automlzero} E.\ Real et al. AutoML-Zero: Evolving Machine
Learning Algorithms From Scratch. \emph{ICML}, 2020.

\bibitem{eureka} Y.\ J.\ Ma et al. Eureka: Human-Level Reward Design via
Coding Large Language Models. \emph{ICLR}, 2024.

\bibitem{voyager} G.\ Wang et al. Voyager: An Open-Ended Embodied Agent
with Large Language Models. \emph{arXiv:2305.16291}, 2023.

\bibitem{sb3} A.\ Raffin et al. Stable-Baselines3: Reliable Reinforcement
Learning Implementations. \emph{JMLR}, 2021.

\bibitem{gymnasium} M.\ Towers et al. Gymnasium: A Standard Interface for
Reinforcement Learning Environments. \emph{arXiv:2407.17032}, 2024.

\bibitem{peft} E.\ J.\ Hu et al. LoRA: Low-Rank Adaptation of Large
Language Models. \emph{ICLR}, 2022.

\end{thebibliography}

\end{document}
